% TOP_PROBLEM: {"id": 30006891, "title": "Hyperkähler SYZ (Lagrangian Fibration) Conjecture", "category_id": 5, "category": {"id": 5, "name": "algebraic_geometry", "display_name": "Algebraic Geometry", "description": "Geometric objects defined by polynomial equations.", "slug": "algebraic-geometry", "order_index": 5, "created_at": "2026-07-31T15:26:25.671Z"}, "status": "open"}
% FIRST_POSTED: 2026-09-27
% !TeX program = lualatex
% ATTEMPT_STATUS: unresolved
\documentclass[11pt]{article}
\usepackage[margin=25mm]{geometry}
\usepackage{fontspec}
\setmainfont{Latin Modern Roman}
\usepackage{amsmath,amssymb,mathtools}
\usepackage{unicode-math}
\setmathfont{Latin Modern Math}
\usepackage{xurl,hyperref}
\hypersetup{colorlinks=true,urlcolor=blue}
\setcounter{secnumdepth}{0}
\setlength{\parindent}{0pt}
\setlength{\parskip}{5pt}
\setlength{\emergencystretch}{3em}
\begin{document}
\section*{\texorpdfstring{Hyperkähler SYZ (Lagrangian Fibration) Conjecture}{Hyperkähler SYZ (Lagrangian Fibration) Conjecture}}\label{30006891}
\subsection{Definitions and mathematical statement}
An irreducible holomorphic symplectic manifold $X$ is compact, simply connected, and K\"ahler, with $H^0(X,\Omega_X^2)$ generated by a nowhere-degenerate holomorphic two-form. Its Beauville--Bogomolov--Fujiki form $q$ is the canonical quadratic form on $H^2(X)$, normalized integrally and characterized up to normalization by $\int_X\alpha^{2m}=c_Xq(\alpha)^m$. A line bundle is nef when its first Chern class lies in the closure of the K\"ahler cone. For every nontrivial nef holomorphic $L$ with $q(c_1(L))=0$, the conjecture is
\[
 \exists a\in{\mathbb{Z}}_{\ge1}:\quad L^{\otimes a}
 \text{ is generated by its global sections}.
\]
This semiampleness assertion is the chosen target. A complete description of the resulting fibration's base or singular fibers is not an additional requirement.
\par\smallskip\textit{Formulation and concept references:} \href{https://epiga.episciences.org/17061/pdf}{[S1]}.

\subsection{Short English statement}
Must every nontrivial nef isotropic line bundle on an irreducible holomorphic symplectic manifold become globally generated after taking a positive power?
\subsection{Research attempt: the K3 surface case and the numerical-dimension barrier}
\textbf{Status.} Semiampleness in every irreducible holomorphic symplectic deformation type is \textbf{UNRESOLVED} by this attempt. We prove the assertion in complex dimension two, derive the precise numerical dimension in arbitrary dimension from the Fujiki relation, and give explicit higher-dimensional examples. The missing general step is production and control of sections; zero Beauville--Bogomolov square is not itself a basepoint-freeness theorem.

\textbf{Source audit.} The published primary source [S1] is stronger than a statement merely about openness of semiampleness. Its Theorem 3.8 gives the result when the pair admits an appropriate deformation with semiample line bundle, and its introduction states the consequence for all currently known deformation classes. This is substantial existing progress. It is not a classification of every possible irreducible holomorphic symplectic manifold, and it does not make that deformation hypothesis automatic for an unknown type. The source also distinguishes semiampleness on the original nef model from a birational fibration. We retain these distinctions rather than treating all uses of the name SYZ as identical.

\subsubsection{1. Nontriviality excludes a hidden flat-line-bundle exception}
For an irreducible holomorphic symplectic manifold $X$, simple connectedness gives $H^1(X,\mathbb C)=0$, and K\"ahler Hodge theory gives $H^1(X,\mathcal O_X)=0$. The exponential sequence therefore makes the first Chern class map
$\operatorname{Pic}(X)\to H^2(X,\mathbb Z)$ injective. Moreover $H^2(X,\mathbb Z)$ has no torsion: the universal coefficient theorem identifies its torsion with torsion arising from $H_1(X,\mathbb Z)$, which vanishes. Thus a nontrivial line bundle has a nonzero integral Chern class, and a torsion line bundle must be trivial.

This matters because on a general compact K\"ahler manifold a non-torsion flat line bundle is nef and numerically trivial but need not have any sections. For instance a degree-zero non-torsion line bundle on an elliptic curve has no nonzero section in any positive power: a section would have a degree-zero effective zero divisor, hence would trivialize the power. Such an example does not satisfy the simple-connectedness hypothesis of the present problem.

For a nonzero nef class $\ell=c_1(L)$ and a K\"ahler class $\omega$, the pairing $q(\ell,\omega)$ is positive when $q(\ell)=0$. The Beauville--Bogomolov form on real $(1,1)$ classes has Lorentz signature, with the K\"ahler cone in one component of the positive cone. A nonzero null vector in its closure pairs positively with every vector in that component. In coordinates $q(x)=x_0^2-\sum_{i>0}x_i^2$, this is the elementary inequality $x_0y_0-x'\cdot y'>0$ for future null $x\ne0$ and future timelike $y$.

\subsubsection{2. The Fujiki identity computes numerical dimension exactly}
Suppose $\dim_\mathbb C X=2m$. Expanding the Fujiki identity at the nef isotropic class gives
\[
 \int_X(\ell+t\omega)^{2m}
 =c_X\bigl(2tq(\ell,\omega)+t^2q(\omega)\bigr)^m.      \tag{1}
\]
The right side has its first nonzero term in degree $t^m$. Comparing coefficients shows
\[
 \int_X\ell^j\omega^{2m-j}=0\quad(j>m),
\]
and
\[
 \binom{2m}{m}\int_X\ell^m\omega^m
 =c_X\,2^m q(\ell,\omega)^m>0.                         \tag{2}
\]
For a nef line bundle its numerical dimension is the largest $j$ for which the left mixed intersection is positive. Therefore
\[
 \nu(L)=m.                                            \tag{3}
\]
In particular $L$ is not big: its top self-intersection vanishes. Replacing $L$ by a positive power changes the positive coefficient in (2) but leaves the vanishing above degree $m$ unchanged.

The usual projective basepoint-free theorem requires a nef bundle whose suitable adjoint difference is nef and big. Here $K_X$ is trivial and every positive multiple of $L$ still has zero top self-intersection, so that theorem cannot be applied by simply setting the adjoint difference equal to a multiple of $L$. The obstruction is an exact failure of the bigness hypothesis, not a question of choosing a larger multiple.

If a power of $L$ is generated, its associated morphism has image dimension $m$. Indeed a globally generated bundle is the pullback of an ample bundle on its image; the largest nonzero mixed self-intersection is the dimension of that image. This agrees with (3). Showing that such a morphism exists is the hard direction, while its expected half-dimensional image is already visible numerically.

\subsubsection{3. Dimension two: Riemann--Roch supplies a moving divisor}
Now let $X$ be a K3 surface, which need not be assumed projective, and let $L$ be nontrivial, nef and satisfy $L^2=0$. The K3 Riemann--Roch formula and Serre duality give
\[
 \chi(L)=2+\tfrac12L^2=2,
 \qquad h^2(L)=h^0(L^{-1}).                            \tag{4}
\]
Fix a K\"ahler class $\omega$. A nonzero effective divisor has positive intersection with $\omega$. Since $L\cdot\omega>0$, $L^{-1}$ cannot be effective. Equation (4) yields $h^0(L)=2+h^1(L)\ge2$.

Write the complete linear system as $|L|=|M|+F$, where $F$ is the fixed effective divisor and $|M|$ has no fixed divisorial component. We use additive notation for line bundles. Since $h^0(M)=h^0(L)\ge2$, $M$ is nonzero and effective. Nefness of $L$ gives
\[
 L\cdot M\ge0,\quad L\cdot F\ge0,
 \quad L\cdot M+L\cdot F=L^2=0.
\]
Thus $L\cdot M=L\cdot F=0$.

Two general divisors in $|M|$ have no common component, so their intersection number is nonnegative: $M^2\ge0$. The Hodge index theorem on a K\"ahler surface says that the intersection form has Lorentz signature on real $(1,1)$ classes. On the orthogonal complement of a nonzero null vector $L$, it is negative semidefinite with radical precisely the line spanned by $L$. Since $M\perp L$ and $M^2\ge0$, it follows that
\[
 M^2=0,\qquad c_1(M)\in\mathbb R_{>0}c_1(L).           \tag{5}
\]
The positive scalar follows by pairing with $\omega$.

There are no base points of $|M|$. Otherwise two sufficiently general members without common components would both pass through a base point and have positive local intersection multiplicity there, contradicting $M^2=0$. The same reasoning includes infinitely near base points only after an actual blowup; already an ordinary base point would contribute positively, so none exist on $X$. Hence $M$ is globally generated.

\subsubsection{4. The moving part determines the original bundle}
The morphism defined by $|M|$ has one-dimensional image. It is nonconstant because $h^0(M)\ge2$ and $M\ne\mathcal O_X$. It cannot have two-dimensional image: two general hyperplanes on a surface image would pull back to divisors of positive intersection, contradicting $M^2=0$.

Take its Stein factorization $f:X\to B$, so $B$ is a smooth projective curve and $f$ has connected fibers. Pullback of holomorphic one-forms is injective for a surjective morphism to a curve: at a regular value, a nonzero form remains nonzero on a tangent vector mapping nontrivially to the base. Since $H^0(X,\Omega_X^1)=0$, $B$ has genus zero and $B\cong\mathbb P^1$.

Let $E=f^*\mathcal O_{\mathbb P^1}(1)$. Then $M=kE$ for some integer $k>0$, $E$ is nef and globally generated, and $E^2=0$. Connected fibers give $f_*\mathcal O_X=\mathcal O_{\mathbb P^1}$, whence
\[
 h^0(E)=2.                                            \tag{6}
\]
We next prove that $E$ is primitive in $\operatorname{Pic}(X)$, avoiding an unproved assertion about multiple fibers. If $E=aD$ with an integer $a\ge2$, then $D$ is nef, nonzero and has square zero. The Riemann--Roch argument in (4) gives $h^0(D)\ge2$. Choose independent sections $s,t$ of $D$. The $a+1$ sections
$s^a,s^{a-1}t,\ldots,t^a$ of $aD$ are linearly independent: a dependence would make the nonconstant meromorphic function $t/s$ satisfy a nonzero polynomial over $\mathbb C$, which is impossible on a connected complex variety. Thus $h^0(E)\ge a+1>2$, contradicting (6).

By (5), $c_1(L)$ lies on the positive real ray generated by $c_1(E)$. The image of $\operatorname{Pic}(X)$ in $H^2(X,\mathbb Z)$ is saturated: by the exponential sequence it is the kernel of the map to the complex vector space $H^2(X,\mathcal O_X)$, whose additive group has no torsion. Thus primitivity of $E$ in the Picard group gives primitivity of $c_1(E)$ in the integral lattice. Both Chern classes are integral, so $c_1(L)=b\,c_1(E)$ for an integer $b>0$. More explicitly, the intersection of an integral lattice with a rational one-dimensional subspace is generated by its primitive lattice vector. Injectivity of the Chern class map then gives $L\cong E^{\otimes b}$.

Since $E$ is the pullback of $\mathcal O_{\mathbb P^1}(1)$, every positive power is globally generated. We have therefore proved more than the requested semiampleness in dimension two:
\[
 L\cong f^*\mathcal O_{\mathbb P^1}(b),\quad b>0,
 \qquad L\text{ itself is globally generated}.        \tag{7}
\]
A general fiber is smooth by generic smoothness in characteristic zero and has genus one by adjunction, since its square and the canonical class are zero. This identifies the expected elliptic or genus-one fibration without requiring a section.

\subsubsection{5. Concrete higher-dimensional examples from an elliptic K3}
Let $S$ be a K3 surface with a genus-one fibration $f:S\to\mathbb P^1$, and let $X=S^{[m]}$ be its Hilbert scheme of length-$m$ zero-dimensional subschemes. We use the standard theorem that $S^{[m]}$ is an irreducible holomorphic symplectic manifold of dimension $2m$.

The Hilbert--Chow morphism followed by the symmetric product of $f$ gives
\[
 F:S^{[m]}\longrightarrow\operatorname{Sym}^mS
 \longrightarrow\operatorname{Sym}^m\mathbb P^1\cong\mathbb P^m.
\]
The last isomorphism sends an unordered $m$-tuple of points to the coefficients of the binary polynomial having those roots, counted with multiplicity. The map is surjective: choose one point in each fiber over any prescribed tuple, with limits handling repeated entries by properness.

Set $A=F^*\mathcal O_{\mathbb P^m}(1)$. Pulling back coordinate sections shows $A$ is globally generated; it is nontrivial since the map is nonconstant and onto. Its top self-intersection is zero because the base has dimension $m<2m$. The Fujiki identity then gives $q(c_1(A))^m=0$, so $q(c_1(A))=0$. As a pullback of an ample bundle, $A$ is nef. Thus these are explicit higher-dimensional examples satisfying every line-bundle hypothesis and the desired conclusion.

Over a divisor of $m$ distinct smooth base points, the fiber is a product of $m$ genus-one curves. This makes the half-dimensional geometry concrete, but is not used to assert a description of every singular fiber. Nor does construction on one Hilbert scheme show that an arbitrary nef isotropic class on an arbitrary deformation is this particular pullback; that passage requires the deformation and monodromy results discussed in [S1].

\subsubsection{6. Why the surface proof does not simply iterate}
On a K3 surface, the Riemann--Roch equality (4) and vanishing of $h^0(L^{-1})$ force a section because the only remaining cohomology group has a negative sign. In dimension at least four, an Euler characteristic is an alternating sum involving several even positive-degree groups as well. Positivity of that integer cannot alone force $h^0>0$: abstractly, it can be accounted for by $h^2,h^4,\ldots$. A vanishing theorem would be needed before turning such numerical information into sections.

The fixed-part argument is also specifically two-dimensional. On a surface, intersections of two effective divisors without common components count nonnegative local intersection multiplicities; zero self-intersection immediately rules out base points. In dimension $2m>2$, the equality $L^{2m}=0$ is consistent with an expected image of dimension $m$, and intersection of two divisors is positive-dimensional. It cannot exclude an arbitrary base locus by the same local point-counting argument.

Finally, (3) proves numerical dimension $m$, not Iitaka dimension $m$. The latter measures the growth of spaces of sections and is not defined by an intersection polynomial alone. Identifying the two, and then obtaining generation by a power, is the abundance-type content of the conjecture. Replacing that step with the word isotropic would be circular.

\subsubsection{7. Verification and the unresolved step}
The coefficient comparison in (1)--(2) is checked by exact symbolic polynomial expansion for small $m$. It verifies the numerical dimension formula and its factors of two, not the existence of sections. The K3 proof separately checks nontriviality, exclusion of the inverse bundle, moving-part base points, the genus of the base, primitivity of the fiber class and the integral proportionality needed in (7).

The explicit Hilbert-scheme examples use the external holomorphic-symplectic structure theorem and an actual pullback of a globally generated bundle. The source's result for known deformation types is attributed as existing work, not reproved or extended to every possible type.

The strongest independent conclusion is (7) in dimension two and the numerical-dimension calculation in every dimension. The first missing step in unrestricted higher dimension is the effective and ultimately semiample realization of the nef isotropic class, or verification of a deformation-to-semiample hypothesis allowing [S1] to apply. The full catalogue assertion remains \textbf{UNRESOLVED} by this attempt.

\subsection{Sources}
\begin{itemize}
\item[S1]Andrey Soldatenkov and Misha Verbitsky, The abundance and SYZ conjectures in families of hyperkähler manifolds, Épijournal 9 (2025), Article 26.
\url{https://epiga.episciences.org/17061/pdf}.
\textit{Formulation source.}
\end{itemize}
\end{document}
